% RESERVED C.3bg. Unregistered draft: no integration before complete independent audits.
\paragraph{C.3bg (Proposed extension, awaiting independent audit): integral ordinary Picard kernels in every positive characteristic.}
Let $p$ be any prime. For every $N\ge3$ we construct an excellent
normal strictly henselian local ring $R$ of dimension $N$ over an
algebraically closed field of characteristic $p$, with residue field
that field, regular puncture $X$, and regular principal punctured
divisor $Y$, such that
\[
 (\mathbb Z/p\mathbb Z)^{\min(3,N-1)}
 \lhook\joinrel\longrightarrow
 \ker\bigl(\operatorname{Pic}(X)_{\rm alg}
             \longrightarrow\operatorname{Pic}(Y)_{\rm alg}\bigr).
\]
Here algebraic equivalence is defined by actual families on the fixed
product with a smooth proper curve. We give a specified generic
construction, followed by a terminating specialization construction
over $\overline{\mathbb F}_p$. Consequently examples exist over every
algebraically closed field of characteristic $p$. For $N=3$ the
assertion specifies a subgroup of rank two; it does not identify the
rank of the entire restriction kernel. No rational failure is asserted.

\emph{1. Specified ordinary quartics in every characteristic.}
Write $z=(z_0,z_1,z_2)$ and let $\mathcal M_4$ be the fifteen monomials
of total degree four. Take thirty algebraically independent variables
$t_M,s_M$ over $\mathbb F_p$, indexed by $M\in\mathcal M_4$, and put
\[
 F_0=z_0z_1z_2(z_0+z_1+z_2),\qquad
 F=F_0+\sum_{M\in\mathcal M_4}t_M M,
 \qquad H=\sum_{M\in\mathcal M_4}s_M M.
\]
Initially take $k=\overline{\mathbb F_p(t_M,s_M)}$ and
$C=V(F)\subset\mathbb P^2_k$.
We prove that these choices have all the required open properties.

For any plane quartic the hypersurface sequence identifies
$H^1(C,\mathcal O_C)$ with
$H^2(\mathbb P^2,\mathcal O(-4))$, with the three positive exponent
triples $a_0=(2,1,1),a_1=(1,2,1),a_2=(1,1,2)$ as a \v{C}ech basis.
Frobenius is given by the coefficients of $z^{pa_i-a_j}$ in
$F^{p-1}$, up to transpose according to the basis convention.
At $F_0$ we have
\[
 F_0^{p-1}=(z_0z_1z_2)^{p-1}(z_0+z_1+z_2)^{p-1}.
\]
For $i=j$, the remaining exponent is $(p-1)e_i$ and its coefficient
is one. For $i\ne j$, the remaining exponent is $pe_i-e_j$ and has
a negative component. Thus the Hasse--Witt matrix at $F_0$ is $I_3$
for every prime, including two and three. Its determinant is a
nonzero polynomial in the quartic coefficients.

The smooth-quartic locus is a nonempty open set. For odd $p$ the
Fermat quartic proves nonemptiness. In characteristic two the quartic
$z_0^3z_1+z_1^3z_2+z_2^3z_0$ is smooth: at a common zero of its
three derivatives a zero coordinate forces all coordinates zero;
otherwise the three displayed monomials are equal and nonzero, so
their sum is nonzero. This excludes every projective singular point.
The condition $F(0,0,1)\ne0$ is also nonempty and open. Irreducibility
of the coefficient space now proves that our generic $C$ is smooth,
ordinary, and avoids $(0,0,1)$. In particular $(z_0,z_1)$ is a
basepoint-free pencil in $H^0(C,K_C)$, where $K_C=\mathcal O_C(1)$.
The ordinary implication uses exactly the Hasse--Witt and
Verschiebung inputs already opened and checked in C.3bf.

\emph{2. A nondegenerate deformation form without a restriction on $p$.}
For a smooth plane quartic the Euler sequence restricted to $C$
shows that
\[
 H^1(C,\mathcal O_C)\longrightarrow
 H^1(C,\mathcal O_C(1))^{\oplus3}
\]
is an isomorphism: its Serre-dual map sends the three standard
basis vectors to $z_0,z_1,z_2$ in $H^0(C,K_C)$. Hence
$H^1(C,T_{\mathbb P^2}|_C)=0$, and the normal sequence gives a
surjection
\[
 \delta:H^0(C,\mathcal O_C(4))\twoheadrightarrow H^1(C,T_C).
\]
Restriction of plane quartics onto $C$ is surjective as well.
Moreover multiplication
\[
 \operatorname{Sym}^2 H^0(C,K_C)
       \xrightarrow{\ \sim\ }H^0(C,K_C^2)
\]
is an isomorphism: both spaces have dimension six and a plane
quartic has no quadratic equation. Serre duality therefore identifies
$H^1(C,T_C)$ with the space of symmetric bilinear forms on
$H^0(C,K_C)$. This identification uses the ordinary symmetric square,
without division by two. The form associated to $\delta(H)$ is
\[
 Q_H(s,u)=\langle\delta(H),su\rangle.
\]
Every symmetric bilinear form occurs, in particular a nondegenerate
one in every characteristic. Nondegeneracy is a nonempty open
condition on $H$, so it holds for the specified generic coefficients.

These arguments also prove existence over $\overline{\mathbb F}_p$.
The locus of pairs $(F,H)$ with the four properties just proved is a
nonempty open set of finite type over $\mathbb F_p$: smoothness,
$F(0,0,1)\ne0$, invertible Hasse--Witt matrix, and nondegenerate
$Q_H$. It has a point over a finite extension of $\mathbb F_p$.
To make this a terminating specification, enumerate finite fields
and pairs of fifteen coefficient tuples, and retain the first pair
satisfying these four tests. Smoothness and the first two explicit
matrix tests are finite polynomial computations; $Q_H$ is computed
by the finite-dimensional Euler/normal/Serre-duality linear maps
just displayed. Existence proves termination; no bound on the
search time is asserted. The resulting data can be base changed
to any algebraically closed field of characteristic $p$.

\emph{3. The ring and the complete source and target Picard calculation.}
Put $m=N-2\ge1$ and $q=\min(3,m+1)$, and set
\[
 \ell=\sum_{i=0}^{q-1}x_i z_i.
\]
In $\operatorname{Tot}(\mathcal O_{\mathbb P^m\times\mathbb P^2}(-1,-1))$
use the same specified contraction as C.3bf, with equation
$F+t\ell H=0$ and coordinates $w_{ij}=tx_i z_j$. Take the reduced
affine image, normalize, localize at its unique vertex and strictly
henselize to obtain $R$; put $f=\sum_{i=0}^{q-1}w_{ii}$.
The graph argument proves integrality because $F$ does not divide
$\ell H$; a nondegenerate $Q_H$ excludes $H$ proportional to $F$.
Smoothness of $C$ proves smoothness along the zero section. The
proper contraction and normalization argument of C.3bf therefore
provides a regular projective resolution $\mathcal Z$ with
\[
 V=\mathbb P^m\times C,\quad
 L=\mathcal O(1)\boxtimes K_C,\quad
 \mathfrak m_R\mathcal O_{\mathcal Z}=I_V,\quad X=\mathcal Z-V.
\]
All these arguments use smoothness, rather than a Fermat derivative
or division by four, and consequently hold also for $p=2$.
The strict transform of $f=0$ is the split line bundle over
$H_0=(\ell=0)\subset V$. The entries of $\ell$ generate $K_C$;
thus $H_0$ is a smooth projective bundle of relative dimension $m-1$
over $C$, and its punctured line bundle is the regular $Y$.
For $m=1$ it is precisely $C$, and $L|_{H_0}=K_C^2$.

As in C.3bf, for every $m\ge1$,
\[
 H^1(V,L)=k^{m+1},\quad H^2(V,L)=0,\qquad
 H^1(V,L^n)=H^2(V,L^n)=0\quad(n\ge2).
\]
The first-neighborhood Picard extension is consequently
\[
 0\longrightarrow U_0:=\mathbb G_a^{m+1}
 \longrightarrow G:=\operatorname{Pic}^0_{V_2/k}
 \longrightarrow J:=\operatorname{Pic}^0_{C/k}\longrightarrow0,
\]
and every later Picard restriction is an isomorphism. The unit
transition maps are surjective after the second stage, so compatible
trivializations have no inverse-limit obstruction.
On the split target the full formal Picard identity component is
\[
 J\times(U_0/U),\qquad
 U=\ell H^1(C,\mathcal O_C)=\langle x_0,\ldots,x_{q-1}\rangle.
\]
Indeed the latter multiplication map is Serre dual to the inclusion
of the independent sections $z_0,\ldots,z_{q-1}$; it has rank $q$.
The sequence $0\to L^{n-1}\to L^n\to L^n|_{H_0}\to0$ proves
$H^1(H_0,L^n|_{H_0})=0$ for $n\ge2$. For $m=1$ the vector quotient
is zero, consistently with $L|_{H_0}=K_C^2$.

The first-neighborhood deformation is
\[
 \kappa=\sum_{i=0}^{q-1}x_i\otimes
             \bigl(z_i\delta(H)\bigr)
 \in H^0(\mathbb P^m,\mathcal O(1))\otimes H^1(C,\mathcal O_C).
\]
The Picard-extension map is the contraction
\[
 c_\kappa:H^0(C,K_C)\longrightarrow U_0,
 \qquad s\longmapsto\sum_{i=0}^{q-1}Q_H(z_i,s)x_i.
\]
Illusie's formula and the actual Poincare transition-unit proof in
C.3bf identify this with the extension class, in every characteristic
\citep[Theorem2.6 and formula(2.7.2), pp.~166--167]{D_Illusie1972Torsors}.
Nondegeneracy of $Q_H$ gives rank $q$ and image $U$.

\emph{4. The actual finite subgroup, including the rank-two case.}
Ordinarity gives an etale Verschiebung
$v:A:=J^{(p)}\to J$ with kernel $K\simeq(\mathbb Z/p)^3$ and
$v^*=0$ on $H^1(\mathcal O)$. The bifunctorial integral extension
classification in \citet[equation(3.8), Section3.2.3, p.~23]{D_BrionIsogenyII}
therefore splits $v^*G$ and gives a homomorphism $b:A\to G$ above $v$.
The Hom--Ext sequence identifies
$\operatorname{Hom}(K,\mathbb G_a)$ with
$\operatorname{Ext}^1(J,\mathbb G_a)$. Thus the $k$-linear span of
$b(K)$ in $U_0$ has dimension $q$, because $c_\kappa$ has rank $q$.
A finite elementary abelian subgroup with $k$-linear span of dimension
$q$ has $\mathbb F_p$-rank at least $q$. Consequently $b(K)$ contains
$(\mathbb Z/p)^q$. When $q=3$, $b|_K$ is injective and this image has
exact rank three. When $q=2$ we make no assertion that the image has
exactly rank two; it may have rank three.

The target restriction of $b$ is $(v,0)$, because its projection
onto $J$ is $v$ and every morphism from the proper connected $A$
to an affine vector group is constant, hence zero when it is a
homomorphism. In particular every element of $b(K)$ has zero
formal target class. We do not assume that $b$ is an immersion
in the rank-two case.

\emph{5. Actual fixed-product families and injective passage to the puncture.}
Take $\alpha\in b(K)(k)$ and choose $\eta\in K(k)$ with
$b(\eta)=\alpha$. Choose a smooth connected proper curve $P\subset A$
through $0$ and $\eta$ by sufficiently ample successive sections.
Pull the normalized actual Poincare bundle from $C\times J$ to
$V\times P$ along $v|_P$. The same mixed Atiyah obstruction as in
C.3bf vanishes because $v^*$ on $H^1(\mathcal O_J)$ is zero.
Hence an actual bundle $\mathcal Q_2$ exists on $V_2\times P$.
Its Picard map differs from $b|_P$ by a constant map to $U_0$;
tensor by the inverse fibre at zero to obtain precisely $b|_P$.
This argument works on the domain $A$ even when $b$ has a kernel.

For $n\ge2$, coherent Kunneth gives
\[
 H^2(V\times P,L^n)=
 H^2(V,L^n)\otimes H^0(P,\mathcal O_P)
 \oplus H^1(V,L^n)\otimes H^1(P,\mathcal O_P)=0.
\]
Choose compatible extensions to all $V_n\times P$. Proper formal
existence on $\mathcal Z_{\widehat R}\times P$, followed by Popescu
finite-presentation descent and smooth lifting over the henselian
pair $(R,\mathfrak m_R^2)$, gives a line bundle on the fixed product
$\mathcal Z\times P$ matching $\mathcal Q_2$. The inverse bundle and
its specified first-jet identification are included in the finite
descent data. These are exactly the formal-existence and approximation
steps checked in C.3bf, C.3az and C.3w, with the same maximal-ideal
identity $\mathfrak m_R\mathcal O_{\mathcal Z}=I_V$.
Individual endpoint Picard lifts are unique from the first jet,
since $H^1(V,L^n)=0$ for $n\ge2$. Proper formal existence and Elkik's
puncture comparison \citep[unnumbered corollary, p.~573]{D_Elkik1973}
make the endpoint at zero actually trivial on $X$. The other endpoint
therefore defines an ordinary algebraically trivial line bundle on $X$.

Distinct $\alpha$ give distinct classes: a line bundle trivial on
$X$ is $\mathcal O_{\mathcal Z}(aV)$; its trivial restriction to $V$
forces $a=0$ by ampleness of $L$, and its first jet then forces
$\alpha=0$. Formal uniqueness and Elkik's injection show that the
assignment is independent of choices and additive. On the target,
the first jet is zero by Step4, all later endpoint lifts are zero,
and proper formal existence plus Elkik for $R/(f)$ proves actual
triviality on $Y$. Thus the whole group $b(K)$ injects into the
ordinary restriction kernel, and its subgroup of rank $q$ proves
the stated assertion. Every step retains its integral coefficients.
